\documentclass[11pt]{article}
\usepackage[margin=1in]{geometry}
\usepackage{amsmath}
\usepackage{booktabs}
\usepackage{longtable}
\usepackage[hidelinks]{hyperref}
\usepackage{caption}
\usepackage{enumitem}

% The default itemize bullet is drawn from a bitmap face, which carries no text
% and shows up as a Type 3 font in the built file. The maths bullet comes from
% an outline face that is embedded anyway.
\setlist[itemize,1]{label={-}}

\hypersetup{
  pdftitle={Bread and Bicycles: What Happens to an Interaction Claim When the
            Second Detail Means Nothing},
  pdfauthor={Michael Podgortsev},
  pdfsubject={Measurement validity in audits of language-model bias: a floor
              check for interaction claims},
  pdfkeywords={language models, bias audit, interaction, intersectionality,
               neutral control, preregistration, reproduction}
}

\title{\textbf{Bread and Bicycles}\\[0.3em]
\large What Happens to an Interaction Claim When the Second Detail Means Nothing}
\author{Michael Podgortsev\\ \small podgortsev@gmail.com}
\date{November 2026\\[6pt]
\normalsize DOI: \href{https://doi.org/10.5281/zenodo.23959679}{10.5281/zenodo.23959679}}

\begin{document}
\maketitle

\begin{abstract}
\noindent
When a person tells a language model two things about themselves rather than
one, does the model treat them worse than the sum of the parts? The published
literature is split: one camp reports compounding disadvantage at an
intersection, a newer camp reports sub-additivity and a threshold past a few
attributes. We put the same stacked signals through five outcome measures on
three open models (Qwen2.5-7B, Llama-3.1-8B, Mistral-7B), with clause count held
constant so that ``two signals'' never also means ``one more sentence''.

Alongside every audited pair we stack pairs of socially \emph{neutral} details
built the same way, in the same slots: cycling to work and baking bread, walking
a dog and learning guitar, making tea and reading history books. Their interaction should be zero. On four of the five instruments it is not.
Trait log probability, the rebuilt head-to-head choice and the log probability
of the correct answer report significant interactions between two meaningless
details, on all three models, on the last of those only within difficulty strata
for Qwen, and in places as large as the social ones; the
mechanism is visible in the data, in that the effect of a clause depends on the
clause beside it. The original head-to-head choice is counted for a different
reason: its interaction contrast is invalid by construction, so it yields no
interaction result at all, and its floor moves as well. Only one instrument, whether the model
solves the person's own arithmetic task, holds a clean zero, on all three models
and within difficulty strata.

On that instrument, one interaction of twenty-four survives correction: on
Llama, in one slot order, two stated disclosures together cost 7.2 accuracy
points \emph{less} than
the sum of their separate costs. It is sub-additive, the opposite of
compounding, and it is the size of the spread the same measurement shows across
six neutral wrappers. Stacking depth produces no shape on any of six ladders. A
credential recovers zero to nine percent of the penalty for stating one's age.
Asked which of two signals moved them, the models cannot say.

We registered before collection what would confirm and what would refute the
hypothesis that signals changing the \emph{text} compound while signals that
merely \emph{state} a fact do not. Form-by-form compounding is not detected on any model or instrument that can
identify it, which gives no support to half of that hypothesis but falls short
of refuting it, since the one clean instrument cannot see small effects; the
other half is unreadable on the instruments that fail their floor check.

Four failing instruments could be a fact about our implementation, so we then
ran a neutral comparison on published work with the authors' own code, under a
second preregistration, reproducing one published number before adding
anything. Two
targets carry a complete result. In a widely cited study reporting that identity interactions are significant in
about a third of regressions, the same statistic returns 0.22 to 0.25 for the
adjectives \emph{hungry} and \emph{sleepy}, against a nominal 0.05; two of its four identity groups sit in that range, and two nominally exceed the
preregistered neutral-range threshold, by 0.002 and 0.004, against a single
shared neutral pair of that size, with no second pair registered; a secondary
randomisation test does not separate them from the neutral pool ($p = 0.33$ and
$0.27$). A 2025 benchmark reports coreference disparities above 20 percent, and in our
rerun of one of its five models a pair of neutral markers reaches 0.350 on
type-2 sentences; eleven of twenty attribute-by-sentence-type cells still sit above every neutral
set, four of them by less than their own interval width. In that same rerun, on point estimates, its
doubly-disadvantaged pattern is largest with no identity marker at all, and
adding one mostly reduces it. On four further targets we could not assemble everything a rerun needs, and we
report the five artifacts such a rerun turns out to depend on.

The contribution we can defend is methodological: a compounding result from an
instrument without a floor check cannot be distinguished from the instrument
responding to more text. Four of the five instruments we built do not return zero on that check, and on
both published instruments we were able to rerun, neutral details move the
statistic as far as many of the social ones. The floor check is released as a
standalone tool.
\end{abstract}

\section{Introduction}

This is the third paper in a programme about what systems cost the people they
were not built around. Paper~1 defined the average-fit loss, the cost of sitting
far from the configuration a system was designed for \cite{p1}. Paper~2 asked
what changes when a language model learns \emph{one} thing about a person, and
measured it through seven channels on the same three models and the same one
hundred people \cite{p2}. This paper asks what happens when there are two.

The question matters because the literature is split, and both camps are
plausible. A large body of work reports that disadvantage compounds at an
intersection: Guo and Caliskan define an \emph{emergent} intersectional bias as
one not predicted from its constituent categories \cite{guo2020}, Lalor et al.
report amplification at intersections across ten downstream tasks
\cite{lalor2022}, ABLEIST reports up to 58-fold ableist harm amplified for
multiply-marginalised candidates \cite{ableist2025}, and Compounding
Disadvantage reports variance rising with each added demographic dimension
\cite{gupta2026}. A smaller and newer body reports the opposite: Fu et al. state
that when two signals are present their effects are ``generally sub-additive
rather than simply additive'' \cite{fu2026}, and Kamruzzaman et al. report an
over-specification threshold past one to three attributes
\cite{kamruzzaman2026}. Paper~2's accumulation method sat in the second camp
with a null.

The literature has not yet directly compared these claims using the same
signals, the same models, the same people and the same outcome measures. That is what this paper
does, and the adjudication turns out to depend on something prior to the
question: whether the instrument can tell two meaningful details from two
meaningless ones.

The paper therefore has three parts. Part~I puts two stacked signals through
five outcome measures on three open models, with a neutral pair stacked the same
way beside every audited pair. Four of the five instruments do not meet that criterion.
Part~II takes up a reasonable objection, that five instruments we built
are five instruments we built, and runs a neutral comparison on published work
with the authors' own code, under a second preregistration that requires one of
their numbers to be reproduced before anything is added. Two targets passed that gate, and in both, neutral details move the statistic
well above zero. Part~III lists what a rerun of this kind turns out to depend on. On four of the
seven targets we could not assemble all of it, and the list is what we now ship
with our own work.

What the reader is asked to take away is one sentence, and it is methodological
rather than about any particular group: an interaction claim produced by an
instrument that has not been shown to return zero for two meaningless details
cannot be told apart from the instrument responding to more text.

\paragraph{Terminology.} We say \emph{signal stacking} rather than
intersectionality. The literature above is about identity categories crossing.
Our catalogue is organised by type of deviation and contains items that are not
identity categories at all: writing as a non-native speaker, imprecise typing
whose cause is unknown, first-time use of a tool. Borrowing a term with a
specific social-theory meaning for that set would be wrong, and the distinction
earns its place because the variable of interest here is the type of signal,
with identity categories as one type among several.

\section{Related work, and what this paper concedes}

The design is not novel and we do not claim it. Hoffmann et al. apply a full
factorial to an LLM hiring audit \cite{hoffmann2026}, and Leyva-Vazquez and
Smarandache run three factorial studies with bootstrap intervals on interaction
magnitudes, documenting a rescue effect in which a prestigious publication
compensates for a low-prestige institution \cite{leyva2026}. Sub-additivity is
not a discovery: Fu et al. published it for identity and opinion \cite{fu2026},
and Liu et al. report that two form corruptions of the task text together cost
more than either alone but, in most settings, less than the sum \cite{liu2026}. The over
specification threshold is published \cite{kamruzzaman2026}. What is new here is
the \emph{contrast} between signal types in one design, the floor check applied
to every instrument, and the finding that most of them do not meet the
preregistered criterion.

Three prior results shape the methodology. First, instrument validity: Wang et
al. show that first-token probabilities mismatch the text an instruction-tuned
model actually emits at rates above 60 percent \cite{wang2024a}, and argue that
the text is the more robust read \cite{wang2024b}. We answer this rather than
assume it: every probability read in this paper records the mass on the answer
tokens, the gate is applied in the analysis, and every reported comparison is a
difference between two conditions in the same slot, so the slot term cancels.
Second, order effects: Yin et al. report that position bias is generally
stronger than gender bias and cannot be removed by averaging or randomising
presentation \cite{yin2025}. That is the external account of why Paper~2's
forced-choice method failed at a position ceiling, and it is why pairs carrying
a stated signal run in both slot orders here. Third, cue instability: Tonneau et
al. report over 14.8 million prompts that both the magnitude and the direction
of a group difference vary with how the group is cued \cite{tonneau2026}. Our
wrapper experiment is the neutral-text version of that warning, and it is the
error bar around every other number we report.

\paragraph{The surrounding literature, in one place.} The question of how
subgroups should be audited at all predates language models: Kearns et al. show
that fairness holding on each attribute separately can fail on their
intersections, and audit subgroups directly \cite{kearns2017}, and Gohar and
Cheng survey the notions and the mitigations that followed \cite{gohar2023}.
Within language models the measurements take several shapes: Boufaied et al.
study intersectional fairness directly \cite{boufaied2026}, Reddy et al. test it
by metamorphic relations \cite{reddy2025}, and Maltbie and Raval report that
false validation of a user's claim varies with the demographics the model
perceives \cite{maltbie2026}. A second group asks what a stated persona does to
the model rather than to its verdict on a person: Gupta et al. find implicit
reasoning deficits in persona-assigned models \cite{guptapersona2023}, and
Occhipinti et al. report that sociodemographic prompting helps some groups and
hurts others \cite{occhipinti2026}. A third concerns signals that are carried by
the text rather than stated, which is our form-against-stated distinction from
the other side: Haq and Saldias separate bias from implicit linguistic signals
from bias from explicit profiles \cite{haq2026}. A fourth measures what
disclosure costs the person who discloses: Kamruzzaman and Kim on disability
disclosure in candidate selection \cite{kamruzzaman2025}, Yildirim on mental
health disclosure to personalised agents \cite{yildirim2026}, Poole-Dayan et al.
on underperformance falling disproportionately on vulnerable users
\cite{poole2024}, and Panda et al.'s disability benchmark across 21 models
\cite{panda2025}. Finally, on the statistics of designs like ours, Pashley et
al. give sample-size determination for $2^K$ audit experiments with a binary
response \cite{pashley2025}, which is the family our channel~A belongs to and
the reason its minimum detectable interaction is printed beside every null.

\paragraph{Controls closest to the floor check.} Three recent works come
closest, and none of them, to our reading, tests an interaction. Vohra and Ravikiran add a placebo to LLM decision
audits in which candidates differ only in a decision-irrelevant detail, a hobby
or a neighbourhood, and find that under a forced-choice probe the placebo reads
as unfair 74 percent of the time, against 5 percent under open-ended scoring
\cite{vohra2026}; their placebo varies one detail and protects a single contrast.
WinoIdentity includes a non-demographic augmentation baseline, adjectives such as
``confused'' and ``relaxed'' in the same slots, used for single-axis gender
disparity rather than for its doubly-disadvantaged claim \cite{winoidentity2025}.
Rennard and Xypolopoulos calibrate intersectional persona effects against
sampling noise and against a random unrelated demographic dimension, and find
that models keep one to two identity dimensions and discard the rest
\cite{rennard2026}; their null is demographic rather than neutral. The pattern our answer-probability channel shows, a second clause adding less
than the first, has the same shape as the sub-additivity Fu et al. report, who
state that ``When both signals are present, their effects are generally
sub-additive rather than simply additive'' \cite{fu2026}; on our instrument that
channel does not pass the floor check, which is the point of Section~4. What none of these does, and what
the floor check adds, is stack two socially neutral details in the structural
positions of the two audited signals and require their interaction to be zero
before any social interaction is read. Magee et al. show why it matters: in their data prompt length and the number of
terms both correlate with sentiment, and they conclude that intersectional bias
``cannot be easily estimated from single-category bias'' \cite{magee2021}.
Separating those two is what a neutral pair in the same slots is for.

\section{Design}

\paragraph{Terms.} Four words recur and are easy to confuse, so they are fixed
here. The \emph{neutral floor} is the interaction estimate for a pair of neutral
details; it should be zero. The \emph{floor check} tests whether it is, and an
instrument whose floor check rejects is \emph{contaminated}. In Part~II, where
the published statistic is not an interaction, the \emph{neutral range} is the
interval that statistic takes for neutral attributes; it is a comparison band,
not a test of zero. The \emph{wrapper spread} is the variation of an estimate
across six neutral opening sentences; it measures sensitivity to wording, not
contamination.


\paragraph{Models.} Qwen2.5-7B-Instruct, Llama-3.1-8B-Instruct and
Mistral-7B-Instruct-v0.3, loaded at 4-bit on a single free-tier T4. Different
developers, different countries, different training data. An effect present in
one model only is a fact about that model, and is reported that way.

\paragraph{Signals.} Four, drawn from the programme's axes catalogue. Two
\textsc{form} signals, which change the text the model must read: writing as a
non-native speaker (F01) and typing errors (F04). Two \textsc{stated} signals,
which leave the task byte-identical and only assert something about the person:
``I use a screen reader for everything on my computer'' (S01) and ``I am
seventy-four and not very confident with computers'' (S08). Four 2x2 pairs:
form-by-form, stated-by-stated, and two cross pairs.

\paragraph{Clause count is held constant.} Every lead-in is exactly two clauses
in every condition. A signal takes a slot that a neutral clause would otherwise
hold. Without this, ``two signals'' also means ``one more sentence'', and the
experiment measures verbosity. Paper~2 measured what that costs: a 40-word
preamble moved Llama's accuracy by ten points regardless of content.

\paragraph{The interaction.} Per item,
\[
d = (\text{both} - A) - (B - \text{none}),
\]
paired within item. Negative is super-additive, the pair costing more than the
sum of its parts, which is the compounding claim. Positive is sub-additive.

\paragraph{The floor.} Alongside the audited pairs we stack pairs of socially
neutral but distinctive details in the same slots under the same arithmetic:
cycling to work with baking bread, walking a dog with learning guitar, making
tea with reading history books. Each was checked at build time against a cue-word
list; a first draft used ``crime novels'' and was rewritten because an instrument
that reads text reads words. Three floors, not one, because a single neutral pair
can itself be unusual.

\paragraph{Instruments.} Five, all reading the same conditions.

\begin{center}
\begin{tabular}{llll}
\toprule
 & what it reads & items & outcome \\
\midrule
A & does the model solve the person's task & 600 tasks & binary \\
A2 & log P(correct answer) on A's prompts & 600 tasks & continuous \\
B & log prob of six favourable words & 100 profiles & continuous \\
C & head-to-head margin against a twin & 200 pairs & continuous \\
C2 & head-to-head, both-cell against each single & 200 pairs & continuous \\
\bottomrule
\end{tabular}
\end{center}

\paragraph{Statistics.} Every estimate is a mean with a percentile bootstrap
interval and a minimum detectable effect, because an interaction is a difference
of differences and carries about twice the variance of a single difference. Each
outcome is tested twice: signed rank, and a sign-flip permutation test on the
mean, which is the quantity the interval describes. On outcomes with a ceiling
the two disagree, and choosing between them after seeing the data would be a
forking path, so the rule is fixed in advance and can only remove claims: a
floor moves if \emph{either} test rejects, and an effect is claimed only if
\emph{both} reject and the mean leans with the signed ranks. Floors are corrected family-wise (Holm) as their own family, one per model and
instrument and never together with the contrasts they protect; audited
contrasts get Benjamini-Hochberg. The task set
was extended from Paper~2's 200 items to 600 after a power analysis measured 20
percent power for the interaction at $n=200$; rows T001 to T200 are Paper~2's
items byte for byte, enforced by a provenance check.

\paragraph{Registered expectations.} Written before the first run.
\emph{Confirmation:} form-by-form super-additive and clearing the wrapper
spread, \emph{and} stated-by-stated indistinguishable from zero, in at least two
of three models on at least two of three channels. \emph{Refutation:}
stated-by-stated super-additive, or form-by-form absent, in two of three models
on channels powered for it. \emph{A null that is still the paper:} no
interaction of either type anywhere, with the floor showing that clause count
explains what movement there is.

\section{The floor result}

The three difficulty strata are fixed cuts on each task's mean answer
probability across all conditions, at 0.8 and 0.2, and a stratum is used only
if it holds at least thirty tasks. The cuts are round probabilities rather than
quantiles of these data, but they were chosen after the first A2 analysis, at
the same time as the stratified check itself, and the limitations say so. The
standalone \emph{Interaction Floor Test} instead splits each floor's items into
thirds by the mean of that floor's own four cells, which is a different
stratification, and reaches the same verdicts.

Table~\ref{tab:floors} is the paper's central result.

\begin{table}[h]
\centering
\caption{Floor verdicts. Three neutral pairs per instrument per model,
corrected within one family per model and instrument: Holm over every floor
test, overall and within three difficulty strata, then Bonferroni across the two
tests. Twelve families, twelve tests in each of the task channels. The fraction
counts the floors that move \emph{in aggregate}; the correction and the verdict
also use the stratum tests, which is why Qwen's A2, quiet in aggregate, is
marked clean in aggregate only. On A2 a further floor moves within a stratum on
Llama and on Mistral without moving in aggregate. Every test behind this table
is listed in Appendix~A.}
\label{tab:floors}
\begin{tabular}{lllll}
\toprule
model & A accuracy & A2 answer log prob & B traits & C2 choice pairs \\
\midrule
qwen & clean & clean in aggregate only & contaminated 3/3 & contaminated 6/6 \\
llama & clean & contaminated 2/3 & contaminated 3/3 & contaminated 6/6 \\
mistral & clean & contaminated 2/3 & contaminated 3/3 & contaminated 5/5 \\
\bottomrule
\end{tabular}
\end{table}

Channel C is absent from the table because its interaction contrast is invalid
by construction: differencing two margins measured against a shared reference
does not isolate the extra cost of a second signal. Its floor is large anyway,
at $+0.50$, $+1.35$ and $+3.56$ on the three models.

On channel B the floors run from $-0.25$ to $+0.59$ and every one of nine is
significant. On Llama they do not even share a sign, so the contamination is not
an offset that a correction could subtract. For scale, Qwen's neutral floor of
$+0.553$ is larger than its stated-by-stated social interaction of $+0.400$. On
C2 every floor is significant on every model, from $+0.35$ to $+6.44$.

Channel A2 fails in a subtler way, and it is the reason the floors are also
tested within strata. On Qwen the three floors are quiet in aggregate
($+0.18$, $+0.13$, $-0.02$ against a minimum detectable effect near $0.17$), yet
floor 1 moves inside the near-ceiling stratum ($-0.03$, decision $p$ $0.038$,
Appendix~A) while far
from the ceiling floors 1 and 2 sit at $+0.51$ and $+0.40$. Two movements in
opposite directions cancel in the average.

\paragraph{It is not run-to-run noise.} The signal-free condition was collected
twice, in separate sessions. Channel A gives the same right-or-wrong outcome on
100.0 percent of tasks on all three models; channel B's largest per-profile
difference between runs is 0.0053, 0.0024 and 0.0020 log probability, against
floor movements of 0.09 to 0.59.

\paragraph{The mechanism.} Re-derived independently of the analysers: the effect
of a clause depends on the clause beside it. Adding ``I cycle to work'' to an
otherwise neutral message moves Qwen's choice margin by $-4.39$ logits; adding
the same clause to the message that already mentions baking bread moves it by
$+2.05$. On Llama, $-0.33$ and $+0.38$. An instrument with that property will report an interaction for any pair whose
clauses interact this way, and on these three instruments every neutral pair
tested did.

\section{What survives: the person's own task}

Channel A is the only instrument whose floor is clean on all three models and in
every stratum. It is also the least sensitive: at $n=600$ it detects an
interaction of fifteen accuracy points with 90 percent power and one of five
points with 13 percent. Every estimate below carries its own minimum detectable
effect.

\paragraph{Main effects.} Only Llama moves. A stated disclosure costs that model
3.3 to 6.5 accuracy points on the person's own arithmetic (screen reader
$-0.050$ and $-0.063$ by slot, age $-0.065$ and $-0.033$, all
$p \le 0.015$ by exact McNemar), and typing errors cost $2.7$ points. On Qwen
and Mistral nothing moves, so on those models there is no main effect for an
interaction to modify, and the interaction nulls describe the instrument rather
than the person's treatment.

\paragraph{Interactions.} One of twenty-four survives correction:
stated-by-stated on Llama, $+0.072$ accuracy (adjusted $p=0.004$, minimum
detectable effect $0.036$), agreeing in sign with the other slot order
($+0.038$, not claimed on its own). It is \emph{sub-additive}: two disclosures
together cost seven points less than the sum of their separate costs.
Form-by-form is indistinguishable from zero on all three models, with the
estimate below the minimum detectable effect in every case.

\section{The other instruments, and why they are still reported}

Channel B is kept in the paper as a demonstrated failure rather than as a
result. It measures a common thing in a common way, the readiness to apply
favourable words to a person, and its own floor check rejects on all three models. Its coherence check fails too: on Mistral and Llama several conditions
move the favourable words and their opposites in the same direction, which is a
scale shift rather than a judgement. Deleting it would leave the methodological
argument resting on an assertion about other people's instruments; keeping it
makes the argument a measurement.

The main effects of these instruments are unaffected by the floor result,
because a single clause replacing a neutral clause has no second clause to
interact with. They replicate Paper~2: disclosing a screen reader lowers the
readiness to apply favourable words on all three models, by up to two log
probability units. Neutral details also lower it, by as much, which is a caution
about reading any single-signal penalty against silence rather than against
another detail.

\section{The rest of the design}

\paragraph{Depth (e2).} Two ladders, form and stated, at depths 0 to 3, clause
count constant. The shape quantity is the second difference
$(D_3-D_2)-(D_2-D_1)$, bootstrapped over tasks. Six ladders, six intervals
covering zero: no shape can be claimed anywhere. Two cells of eighteen move at
all, both on Qwen's stated ladder, where two or three signals cost about three
accuracy points against none while one costs nothing measurable. That is a
threshold rather than a slope, and it sits inside the wrapper spread below.

\paragraph{Wrappers (e1w).} Every e1 interaction measured six times, changing
only the neutral opening sentence. The spread across wrappers is three to
fifteen times the mean, and eight of nine interactions change sign between one
neutral opening and another. Nothing clears its own wrapper spread. The one
surviving channel A interaction, $+0.072$, is the size of that contrast's
wrapper spread, $0.080$, which is why it is reported as a direction with a
stated limit.

\paragraph{Legibility (e2b).} Liu et al. stack two form corruptions of the task text
and report a combined drop larger than either factor alone but, in most settings,
smaller than the sum, with occasional super-additive cases \cite{liu2026}. Both of
their factors corrupt the text
the model must read, so we ran the same form-by-form pair in two arms: the
signal in the person's message with the task byte-identical, and the styling
applied to the task as well. Neither arm is super-additive on any model. On
Llama the corrupted-task arm is significantly \emph{sub}-additive ($+0.053$,
$p \le 0.0005$ on both tests), which is the direction Liu et al. report in their
own design, and the two arms differ by $+0.063$. With the task byte-identical the
interaction is null on all three models. So their sub-additive direction appears
on one model of three when the task is corrupted, and nothing appears when only
the person's message carries the form signal, with minimum detectable effects of
0.025 to 0.035 stating what this run could have seen.

\paragraph{Offsetting (e3).} Does a credential buy back a disclosure? On Llama's
three usable numeric measures, stating fifteen years of experience recovers zero
to nine percent of the penalty for stating one's age, and in two of six cells
the recovery is negative. On the choice channel the answer is slot-dependent:
51 percent recovery on Llama when the credential comes first, nothing on any
model when it comes second. The credential is not neutral on its own: on Qwen a
candidate who states fifteen years of experience is chosen \emph{less} often
than one who says nothing distinctive.

\paragraph{Attribution (e5).} Replay the exchange, then ask the model how much
each of the two signals changed its answer. Llama answers zero to all 840
ratings: its self-report is a constant and says nothing, the second time a model
in this programme has failed that way. Qwen ranks the two signals correctly in
16 of 22 informative items, which does not separate from chance after
correction. One of six testable item-level correlations is significant
($\rho=-0.229$). Qwen attributes 9.65 of 10 to cues the person never mentioned, against 3.15 for
what the person actually said. Mistral ranks the two signals correctly in 2 of 9
informative items, 22 percent, which does not separate from guessing.

\paragraph{Refusal (e6).} Between zero and two refusals per 10,200 generations,
at every signal count, on every model. Refusal does not stack on this task.

% Part II of the merged paper: the floor check on published instruments.
% Draft of 2026-10-06. Inputs after "The rest of the design", before "Discussion".
% Every number here comes from a RESULT document in experiments/r*/results/.

\section{A neutral arm on two rerunnable published instruments}

The floor check above is applied to instruments we built. A reasonable objection
is that four of the five instruments not meeting the preregistered
neutral-floor criterion may be a fact about our implementation rather than
about the field. This section answers it by running a neutral comparison on published work with
the authors' own code.

\paragraph{What these reruns rest on.} Every result in Part~II and Part~III
rests on public materials alone: released code, released data, released prompts,
released models. Nothing reported here depends on anything supplied privately,
and where an artifact was not in a release we worked from what was, rather than
reporting a number that rests on something a later reader could not obtain. That
is deliberate, because the question is what a release on its own supports.

The procedure is fixed in a preregistration separate from the one governing
Part~I, with thirteen dated addenda, and it has two steps. \textbf{Step~0, fidelity:} reproduce one published number with the authors' released code and data, under a
rule written before the output is read. Where our rebuild of an instrument does not match theirs at step~0, no neutral arm
is built on it, because an instrument that does not match theirs cannot be used
to read their claim. \textbf{Step~1, the neutral arm:} add attributes that are socially
neutral but otherwise built the same way, in the same slots, and read the
published claim against them.

Seven targets were selected from thirty-one screened papers. Two passed step~0
and carry a complete neutral-arm result; they are reported here. One passed on the
authors' data and came in outside tolerance on our own regeneration. On four we
could not assemble everything a rerun needs, and that is
Section~\ref{sec:materials}.

\subsection{Kirk et al.: what the statistic returns for neutral adjectives}

Kirk et al. report that in a regression of predicted occupation on identity
terms, ``Interactions are significant in approximately one third of
regressions''~\cite{kirk2021}. We regenerated their completions with GPT-2
124M, 38 templates and 7{,}000 completions per template, extracted job titles
with the same CoreNLP version, and ran their regression functions unchanged.

Step~0 passed: the mean share of significant woman-by-religion terms is 0.242
against their 0.203, inside the centred interval of their own released data,
$[0.167, 0.245]$.

The neutral arm replaces the identity descriptors with neutral adjectives
(hungry, sleepy, thirsty, busy) in the same positions, and reads the share
against the groups that can be formed from that pool.

\begin{center}
\begin{tabular}{lccccl}
\toprule
group & terms & social & neutral & neutral interval & verdict \\
\midrule
religion & 5 & 0.242 & 0.247 & $[0.199, 0.290]$ & inside the neutral range \\
ethnicity & 4 & 0.224 & 0.221 & $[0.178, 0.270]$ & inside the neutral range \\
sexuality & 2 & 0.289 & 0.216 & $[0.147, 0.287]$ & above the neutral range \\
political & 2 & 0.291 & 0.216 & $[0.147, 0.287]$ & above the neutral range \\
\bottomrule
\end{tabular}
\end{center}

The reading is in the two right-hand columns. For neutral adjectives such as
\emph{hungry} and \emph{sleepy} this statistic returns 0.22 to 0.25, where
nominal chance is 0.05. For religion and ethnicity the identity share sits in
the same range. For sexuality and political affiliation the identity share sits above the
neutral interval, by 0.002 and 0.004. \textbf{The preregistered rule that decides this is the one stated in the
registration: a group counts as above the neutral range when its mean lies above
the neutral interval at every matched group size.} Both do, so both are above the range under that rule, by a margin of
two to four thousandths and against a single neutral pair of the same size.

A second test was registered as secondary and decides nothing: a randomisation
over every size-two group formable from the neutral pool gives $p = 0.33$ and
$0.27$, at a resolution of 0.067, so it does not separate these two groups from
the groups the neutral pool can form. It is reported because it was registered,
and it should not be read as support for the result, only as a statement that
it cannot distinguish them either way at that resolution.

One possible source of the number is visible in the data. Many per-job regressions have an empty
category-by-gender cell: 41 of 64 for the neutral religion group, 41 of 67 for the identity one, and 38 of 63 in the released data. Logistic regression
with an empty cell separates, and separated terms return extreme $p$-values.
We read this as a property of the regression rather than of the words in the
slot, which would explain why the share is of similar size for the neutral
adjectives. That reading is our own, offered as a possible mechanism; the
reproduced numbers above stand apart from it and do not depend on it.

\subsection{WinoIdentity: where the neutral markers fall}

WinoIdentity augments WinoBias referents with identity markers and reports
finding ``coreference confidence disparities higher than 20\% in 7 out of 10
demographic attributes''~\cite{winoidentity2025}. We scored one of their five models,
Mistral-7B-Instruct-v0.2 at 4-bit, with their prompt format. Step~0 passed:
Spearman 0.905 against their 22 published values, mean absolute difference
0.029.

The neutral arm puts neutral markers in the marker slot. Of twenty
attribute-by-sentence-type cells, eleven sit above every neutral set and nine
are not distinguishable from one. The number to read beside them is the neutral pair itself: \emph{early-rising} and \emph{late-rising} reach a disparity of
\textbf{0.350} on type~2 sentences, and a three-marker neutral set reaches
0.340. On type~2 sentences, a neutral marker produces a disparity above 20 percent on its own. Four of the eleven cells above the neutral range clear their neutral set by
less than the width of their own bootstrap interval.

The second measurement concerns the pattern read as exacerbated bias towards
doubly-disadvantaged identities. Writing $D$
for the mean feminine-minus-masculine coreference confidence on male-dominated
occupations, more negative meaning the feminine version loses more:

\begin{center}
\begin{tabular}{lc}
\toprule
marker & $D$ \\
\midrule
\textbf{no marker at all} & $\mathbf{-0.721}$ \\
neutral markers & $-0.489$ to $-0.676$ \\
able-bodied & $-0.677$ \\
non-English-speaking & $-0.718$ \\
white & $-0.264$ \\
Asian & $-0.241$ \\
neurodivergent & $-0.245$ \\
transgender & $-0.179$ \\
\bottomrule
\end{tabular}
\end{center}

What the table shows, as measured: the feminine disadvantage is largest before
any marker is added, and adding one, identity or neutral, mostly reduces it. By
the preregistered rule, eight of ten attributes fall in the registered category
of a gender pattern carried through any marker. For the two remaining
attributes the most negative marker for disability is \emph{able-bodied}, and
non-English-speaking is no more negative than no marker at all.

How we read it, which is separate from the measurement: in this model and this
evaluation setting the observed pattern is consistent with the gender asymmetry
present in the underlying sentences. That is one reading of these numbers and
not the only one, and further tests would be needed to determine how much
additional variation is attributable to identity markers.

\subsection{What these two reruns do not establish}

Both are single-model. R1 uses GPT-2 124M, as their released regressions do, and
says nothing about their larger models. R2 uses one of five models at 4-bit. In
R1 our social data are regenerated, and fidelity passed on religion, the group
whose released regressions reproduce exactly; the two groups that nominally exceed the neutral range share one neutral pair,
and a second pair of that size was not registered. In R2 the
preregistration asked for a paired bootstrap interval on $D$ which the analyser
does not yet print, and the rule decides on point values as registered. Neither rerun shows that a published claim is false, and neither is offered as
one. What both show is that in these reruns the statistics alone did not clearly
separate the reported effect from responses to a detail carrying no social
content, which is why a neutral arm is worth adding before an effect is sized.


% Part III of the merged paper: what could not be rerun, and why.
% Draft of 2026-10-06. Inputs after part2-published.tex.
% The R6 subsection has a slot for reading B, which is still collecting.

\section{What a rerun needs}
\label{sec:materials}

The seven targets were chosen because each released prompts, or data and code.
On four of them we could not assemble everything a rerun needs. This section
lists what was missing in each case, because it is the cost of Part~II and
because the same list is what we now ship with our own work.

The claim here is deliberately small. Reproducibility has been measured at
scale, on samples of hundreds and thousands of papers, and this is not that. Our
seven targets were selected for a different purpose, and seven cases support no
rate. What they can offer is a concrete list of the artifacts a rerun of this kind
turns out to depend on, which is more useful than a percentage.

\subsection{Five artifacts a rerun of this kind depends on}

\paragraph{Classifier weights.} The LABE metric is computed with three trained classifiers~\cite{wan2025}. The released folder carries configuration, scheduler and trainer state for each.
The weight files are not part of the public release, so we retrained all three
with their own script, reaching validation accuracies of 0.944, 0.933 and 0.941
against the reported 0.933, 0.933 and 0.917. Step~0 then came in at a median absolute relative difference of 0.286 against a
registered tolerance of 0.25, and that difference cannot be separated from
retraining noise: three classifiers of near-identical accuracy produce variances
on the same generations differing by up to a factor of four. The reported ordering, with the intersectional variance largest, appears in 14
of 18 per-classifier rows. That ordering may also be expected from an eight-cell
variance containing both the gender and the race spread, which is one reason the
neutral arm was designed for this target. With the original weights the
question the neutral arm was built for could be settled; without them it stays
open.

\paragraph{The input to a preprocessing tool.} Wilson and Caliskan code resumes and job descriptions to occupations with a
public coding tool~\cite{wilson2024}. A rerun needs to know which text was sent
to that tool, and we could not determine it. We coded 4{,}859 distinct texts
through that tool's API and tried seven readings of the resume side, three
registered before the first attempt and four after it failed and before they
were computed. None of the seven reproduces the published counts, and the difference has the
same shape every time, concentrated in three of the nine occupations. The
consistency of that shape is what suggests a single preprocessing choice we did
not recover, rather than scattered noise. The job-description side reproduces under the first reading, 8 of 9 counts
within 20 percent, so what we were missing is specifically the resume-side
input.

\paragraph{A runtime version.} Ghanbari Haez and Dragoni generated their answers through a local model
server~\cite{ghanbari2025}. Which version of it produced the released answers
is something we were unable to establish, and generation is sensitive to it. Our
regeneration missed the registered fidelity tolerance by 0.0006. A serial rerun,
registered as a diagnostic before it was run, did not close the gap: their step
is 0.8583, our parallel run 0.8680, our serial run 0.8721. Request batching is
therefore not the explanation, and no conclusion about their claim is drawn. What our own generation does show is
reported separately and is a statement about our answers only.

\paragraph{How the comparison groups were formed.} ABLEIST reports compounding
harm across stacked attributes~\cite{ableist2025}. Its Table~6 caption leaves
several groupings open, and all three of our registered readings came in at best
6 of 12 numbers per model against a rule of 10. One mixed reading, which we
reconstructed by search over the readings the caption leaves open, reproduces
all 24 numbers exactly on both models: the three single-and-double columns are
unrestricted, covering every gender and both castes, while the triple-stacked
column is restricted to the marginalised values. Since that reading reproduces every published number to the second decimal on
two independent models, it is very likely the computation behind the table,
though we cannot establish that from the released materials alone. It also
indicates that the composition of the comparison groups may influence the
observed increase, and that this is worth considering when interpreting that
row. The reconstruction did not satisfy the preregistered reproduction criterion, and
the reconstructed interpretation is reported alongside that outcome.

\paragraph{The frame a fine-tuned model expects.} The same work releases a
fine-tuned labeller. Running it needs the instruction string it was trained
with, which we inferred from the released inference code. The labeller could not
in the end be run at the scale our check needed, which is a compute limit on our
side and is reported as one.

\subsection{Beck et al.: how close the reproduction comes}

Beck et al. report that combining all sociodemographic attributes produces the
largest change in predictions for most of their datasets~\cite{beck2024}. Their
code, their sampling notebook and six small open models make this the one target
that is fully runnable, and we ran their pipeline unchanged over six models and
seven conditions.

The registered reading required the printed winner to be ours and our combined
share to be within five percentage points, on at least four of six models.
Three of the six meet both conditions. On a fourth the combination is also the
most influential condition, so the direction of the claim holds on four of the
six, but our share for that model is 13.02 against 7.00 and so misses the
five-point band. The two remaining models put a single attribute first instead:
one changes between 0.0 and 0.1 percent of labels in every condition, so its
ordering turns on a single label, and the other is inside the five-point band on
size while ordering gender first by 2.1 points.

The preregistration names a fallback reading, the same comparison averaged over
all three of their prompt formats, to be used if the first fails. \textbf{It is not reported here, because its collection was not finished in
time.} The first
format is complete and is what the paragraph above rests on; the other two are
partial. No number from them appears in this paper, and the fallback reading
stays open rather than being decided on part of its data.

What the first format does carry is a check on this work rather than on theirs.
The collection was lost to a deleted working directory in October and rebuilt
from nothing: a new environment at a new path from the versions the authors pin,
the source data downloaded again, the sample rebuilt from their notebook. All
thirty-six cells of the recollection reproduce the first to the second decimal,
and the verdict is unchanged. The distilled predictions are in the repository
with a manifest of checksums, so the table above can be recomputed without the
raw collection.

\subsection{What this section is for}

It is not a reproducibility audit and it is not a judgement on any of these
studies, none of which is shown here to be wrong. Of the five artifact types listed above, two could be reconstructed or inferred
from what was released, with different levels of confidence and neither fully
verified: the grouping behind the published cells, which reproduces every
published number but was found by search after the registered readings did not
match and which did not meet the preregistered criterion, and the frame a
fine-tuned model expects, which we inferred from the released inference code but
could not verify at the scale our check needed. The other three could not, which
is worth knowing in both directions. The point is forward-looking and applies first to us: a
rerun of an interaction claim turns out to depend on a short and predictable
list, namely the trained weights behind any learned metric, the exact input to
any preprocessing tool, the runtime version used for generation, the grouping
behind each published cell, and the frame a fine-tuned model expects. Part~I and Part~II of this paper ship each of these that applies to them, and
that is the habit we would encourage.


\section{Discussion}

\paragraph{The registered expectations, answered.} Form-by-form compounding is not detected on any model or instrument that can
identify an interaction. That provides no support for the half of the hypothesis
that says text-changing signals compound, and it is weaker than a refutation:
the registration asked for absence on channels powered for it, and the one
instrument with a clean floor detects fifteen accuracy points with 90 percent
power and five points with 13 percent, so a small compounding effect would not
have been seen. The other half, that stated signals do not compound, is not
contradicted, but it is not confirmed either. Only one instrument can identify
an interaction, it is powered for large effects only, and on one model it finds
a stated-by-stated interaction that is non-zero and sub-additive, which does not
meet the registered criterion of indistinguishable from zero. On the other four
instruments the question cannot be answered. The outcome that matches the data is the third registered one: no compounding
is detected on the instruments that pass the floor check, and the floor check
indicates that, for several instruments in this
study, the observed movement cannot be cleanly separated from responses to
additional text.

\paragraph{What the split in the literature might be.} Our data are consistent
with a simpler account than either camp's. On channel A2, on every model, the
interaction takes the sign opposite to the two single effects: neutral clauses
cost and two of them cost less than the sum; stated signals raise the answer
probability on two models and two of them raise it less than the sum. Of twelve
interactions that pass both tests on that channel, seven are a benefit that does
not add, two are a cost that does not add, three have singles pointing opposite
ways, and none is compounding harm. A second clause adds less than the first,
whichever way the clauses push. An instrument that moves for any pair of clauses would produce a pattern of this
shape whatever the clauses said, with the sign set by the direction of its single
effects. That is a reason to add a neutral arm before sizing an interaction, rather than
a reason to doubt any particular result.

None of this is evidence for a general sub-additive effect. The one interaction
on a clean instrument is one model, one instrument and one slot-order pair, and
it is the size of its own wrapper spread. The A2 pattern is offered as the shape
a text-sensitive instrument would produce, not as a finding about how signals
combine.

\paragraph{What the published reruns add.} A reasonable objection to Part~I is
that five instruments we built are five instruments we built. Part~II answers it in the most direct way available, by running a neutral
comparison on other people's published instruments with their own code, after
reproducing one of their numbers first.
In both targets that passed that gate, neutral details move the statistic well
above zero: a statistic whose neutral value is five times nominal chance, and a
threshold that a neutral marker reaches. In the second, in our rerun of one of its five models and on point estimates,
the pattern read as compounding at an intersection is largest before any
identity marker is added. We do not claim either
published result is false. In these reruns the statistics alone did not clearly separate the reported
effects from responses to socially neutral details. Neutral-condition
comparisons may therefore help clarify the interpretation of such results.

\paragraph{What a failed floor check does not show.} It does not show that a
social interaction is an artifact. It shows that the instrument has a non-zero
baseline interaction, which has to be accounted for before the social one can be
sized; on channel B even that is not possible, because its floors do not share a
sign.

\paragraph{What a reader should take from this.} Before reporting that two
characteristics compound, stack two socially meaningless details the same way
and show that nothing happens. On four of the five instruments here the neutral pair moves the outcome, and on
both published instruments we were able to rerun it moves the outcome too.
The check costs one extra condition set, and it adds a useful diagnostic for
assessing whether an observed interaction reflects the target phenomenon or
sensitivity to additional text.

\section{Limitations}

\begin{itemize}[nosep]
\item Three open models at 4-bit, one task family for the behavioural channel,
four signals, two of each type. Nothing here speaks to proprietary models,
longer conversations, or signals outside this catalogue.
\item Single-wrapper estimates carry the spread measured in e1w as unreported
uncertainty. That includes every number in this paper except the wrapper
experiment itself.
\item Channel A is not powered for small interactions: five accuracy points are
invisible at this $n$, and the paper says so beside every null rather than
calling them zero.
\item The one surviving interaction is one model, one channel, one slot order
pair, and the size of its own wrapper spread.
\item The stratified floor check was added after seeing the first A2 analysis.
It can only remove claims, and it was validated on synthetic data built to
reproduce the failure before being applied, but it was not registered in
advance. The standalone \emph{Interaction Floor Test}, which stratifies
differently, reaches the same verdicts.
\item Mistral's letter-probability reads are heavily gated (32 percent of C2
reads carry less than half the mass on the two letters), so its numbers on that
channel rest on as few as twelve items and are reported with their $n$.
\item Part~II rests on two reruns, not on a survey. Both are single-model: one
uses the 124M model whose regressions its authors released, the other one of
five models at 4-bit. Neither says anything about the larger models in either
paper, and two cases cannot show how common the problem is.
\item In the first rerun the two groups that nominally exceed the neutral range share
a single neutral pair, and a second neutral pair of that size was not registered. In the
second, the preregistration asked for a paired bootstrap interval on the
doubly-disadvantaged statistic which the analyser does not yet print; the rule decides on point values, as registered, and four of the eleven cells
above the neutral range clear their neutral set by less than the width of their
own bootstrap interval.
\item Part~III lists five artifacts a rerun of this kind depends on, drawn from
seven targets. Seven cases
selected for a different purpose support no rate, and nothing there should be
read as one.
\end{itemize}

\section{Reproduction}

Every number in this paper is produced by scripts in the accompanying
repository from the collected CSVs, with no manual steps. The analysis is
reproduced by running the analysers in \texttt{experiments/*/scripts/} against
\texttt{experiments/*/outputs/}; the offline validation suite
(\mbox{\path{validate_*.py}}, \path{reverify_findings.py},
\path{check_provenance.py}) runs without a GPU and checks the statistics on
synthetic data with known answers, re-derives every headline number
independently of the analysers, and verifies that the items inherited from
Paper~2 are unchanged. The floor check is also released as a standalone tool,
\emph{Interaction Floor Test}, which reaches the same verdicts on these data
through a different stratification.

Part~I and Part~II are governed by two separate preregistrations,
\texttt{PREREGISTRATION.md} and \texttt{PREREGISTRATION-PART2.md}, the second
with thirteen dated addenda. They are published as written and are not merged
retroactively, which is the point of having them: every reading tried is in
them, including the ones that failed, and the order in which they were tried was
fixed before the data were seen. Each carries the date and the commit hash of the commit that first contained
it, in a private working repository that is not published, because its history
also holds material that is not ours to publish. Both documents are archived with their own DOI, Part~I at
\href{https://doi.org/10.5281/zenodo.23959847}{10.5281/zenodo.23959847} and
Parts~II and III at \href{https://doi.org/10.5281/zenodo.23960058}{10.5281/zenodo.23960058},
which makes them immutable from the date of that deposit onward. Neither that deposit nor a commit hash proves the earlier date on which
each was written: that much is asserted, supported by a history we can show on
request, and the paper claims no more than that. Each target's reruns live in
\texttt{experiments/r*/}, beside the \texttt{e*} experiments of Part~I, with the
same shape: scripts, outputs, and a result document stating what cannot be
claimed.

\section*{Acknowledgements}

Part~II of this paper exists only because seven research groups released their
work in a form that someone else could run. We reran instruments built by Kirk
et al.~\cite{kirk2021}, by the authors of WinoIdentity~\cite{winoidentity2025},
by Wan and Chang~\cite{wan2025}, by Wilson and Caliskan~\cite{wilson2024}, by
Ghanbari Haez and Dragoni~\cite{ghanbari2025}, by the authors of
ABLEIST~\cite{ableist2025}, and by Beck et al.~\cite{beck2024}. In every case we
used their code, their prompts, their data or their released model, unchanged
wherever we could. Releasing those materials is what made an independent check possible at all.

Section~\ref{sec:materials} lists the artifacts a rerun of this kind depends on.
It is a list of what is easy to leave out, written by an author who left several
of them out of earlier releases of this work, and it is not a comment on the
generosity of the releases above.

\appendix
\section{Every floor test behind Table~\ref{tab:floors}}

Ninety-eight rows, each carrying both tests: each of three neutral pairs
on each instrument and model, in aggregate, and on the two task channels within
each difficulty stratum as well. One family per model and instrument, Holm
within a family for each test and Bonferroni across the two, exactly as
Section~4 describes. The two Holm columns are the per-test values \emph{before}
the Bonferroni step; the decision is made on the last column, which is twice the smaller of those
two, capped at 1, and a floor moves when that value is below 0.05. The rows are written out by
\path{analyse_e1_floors.py} itself, after the correction that produced
Table~\ref{tab:floors}, so the two cannot drift apart; the same rows are in
\path{appendix_a_floor_tests.csv} in the repository.

\begin{center}
\footnotesize
\begin{longtable}{llllrrrrrrr}
\toprule
model & ch. & neutral pair & scope & $n$ & interaction & \multicolumn{2}{c}{raw $p$} & \multicolumn{2}{c}{Holm $p$} & decision \\
& & & & & & rank & mean & rank & mean & $p$ \\
\midrule
\endfirsthead
\toprule
model & ch. & neutral pair & scope & $n$ & interaction & \multicolumn{2}{c}{raw $p$} & \multicolumn{2}{c}{Holm $p$} & decision \\
\midrule
\endhead
\bottomrule
\endfoot
llama & A & cycling, bread & aggregate & 600 & +0.0100 & 0.508 & 0.587 & 1.000 & 1.000 & 1.000 \\
llama & A & dog, guitar & aggregate & 600 & -0.0133 & 0.404 & 0.468 & 1.000 & 1.000 & 1.000 \\
llama & A & tea, history & aggregate & 600 & +0.0150 & 0.401 & 0.461 & 1.000 & 1.000 & 1.000 \\
llama & A & cycling, bread & near ceiling & 160 & +0.0125 & 0.727 & 0.728 & 1.000 & 1.000 & 1.000 \\
llama & A & cycling, bread & middle & 98 & +0.0000 & 1.000 & 1.000 & 1.000 & 1.000 & 1.000 \\
llama & A & cycling, bread & far from ceiling & 342 & +0.0117 & 0.433 & 0.556 & 1.000 & 1.000 & 1.000 \\
llama & A & dog, guitar & near ceiling & 177 & -0.0056 & 1.000 & 1.000 & 1.000 & 1.000 & 1.000 \\
llama & A & tea, history & near ceiling & 177 & +0.0565 & 0.025 & 0.042 & 0.279 & 0.460 & 0.558 \\
llama & A & dog, guitar & middle & 96 & -0.0104 & 0.906 & 1.000 & 1.000 & 1.000 & 1.000 \\
llama & A & tea, history & middle & 96 & +0.1250 & 0.159 & 0.199 & 1.000 & 1.000 & 1.000 \\
llama & A & dog, guitar & far from ceiling & 327 & -0.0183 & 0.109 & 0.107 & 1.000 & 1.000 & 1.000 \\
llama & A & tea, history & far from ceiling & 327 & -0.0398 & 0.007 & 0.011 & 0.081 & 0.134 & 0.161 \\
llama & A2 & cycling, bread & aggregate & 600 & +0.1401 & 0.009 & 0.008 & 0.027 & 0.041 & 0.054 \\
llama & A2 & dog, guitar & aggregate & 600 & +0.1353 & $<$0.001 & $<$0.001 & $<$0.001 & 0.003 & $<$0.001\,$\star$ \\
llama & A2 & tea, history & aggregate & 600 & +0.2139 & $<$0.001 & $<$0.001 & $<$0.001 & 0.001 & $<$0.001\,$\star$ \\
llama & A2 & cycling, bread & near ceiling & 59 & +0.0213 & $<$0.001 & 0.015 & $<$0.001 & 0.060 & $<$0.001\,$\star$ \\
llama & A2 & dog, guitar & near ceiling & 59 & +0.0270 & $<$0.001 & 0.003 & $<$0.001 & 0.017 & $<$0.001\,$\star$ \\
llama & A2 & tea, history & near ceiling & 59 & +0.0713 & $<$0.001 & $<$0.001 & $<$0.001 & 0.001 & $<$0.001\,$\star$ \\
llama & A2 & cycling, bread & middle & 123 & -0.1567 & 0.034 & 0.033 & 0.067 & 0.098 & 0.134 \\
llama & A2 & dog, guitar & middle & 123 & -0.0174 & 0.080 & 0.783 & 0.080 & 0.783 & 0.160 \\
llama & A2 & tea, history & middle & 123 & +0.1248 & $<$0.001 & 0.103 & 0.001 & 0.207 & 0.002\,$\star$ \\
llama & A2 & cycling, bread & far from ceiling & 418 & +0.2442 & $<$0.001 & $<$0.001 & 0.003 & 0.006 & 0.005\,$\star$ \\
llama & A2 & dog, guitar & far from ceiling & 418 & +0.1955 & $<$0.001 & $<$0.001 & $<$0.001 & 0.002 & $<$0.001\,$\star$ \\
llama & A2 & tea, history & far from ceiling & 418 & +0.2603 & $<$0.001 & $<$0.001 & $<$0.001 & 0.001 & $<$0.001\,$\star$ \\
llama & B & cycling, bread & aggregate & 100 & -0.0908 & $<$0.001 & $<$0.001 & $<$0.001 & $<$0.001 & $<$0.001\,$\star$ \\
llama & B & dog, guitar & aggregate & 100 & -0.2535 & $<$0.001 & $<$0.001 & $<$0.001 & $<$0.001 & $<$0.001\,$\star$ \\
llama & B & tea, history & aggregate & 100 & +0.2080 & $<$0.001 & $<$0.001 & $<$0.001 & $<$0.001 & $<$0.001\,$\star$ \\
llama & C2 & cycling, bread & aggregate, vs a & 200 & +0.6303 & $<$0.001 & $<$0.001 & $<$0.001 & $<$0.001 & $<$0.001\,$\star$ \\
llama & C2 & cycling, bread & aggregate, vs b & 200 & +0.7013 & $<$0.001 & $<$0.001 & $<$0.001 & $<$0.001 & $<$0.001\,$\star$ \\
llama & C2 & dog, guitar & aggregate, vs a & 200 & +0.7966 & $<$0.001 & $<$0.001 & $<$0.001 & $<$0.001 & $<$0.001\,$\star$ \\
llama & C2 & dog, guitar & aggregate, vs b & 200 & +0.4688 & $<$0.001 & $<$0.001 & $<$0.001 & $<$0.001 & $<$0.001\,$\star$ \\
llama & C2 & tea, history & aggregate, vs a & 200 & +0.3509 & $<$0.001 & $<$0.001 & $<$0.001 & $<$0.001 & $<$0.001\,$\star$ \\
llama & C2 & tea, history & aggregate, vs b & 200 & +0.3822 & $<$0.001 & $<$0.001 & $<$0.001 & $<$0.001 & $<$0.001\,$\star$ \\
mistral & A & cycling, bread & aggregate & 600 & +0.0217 & 0.221 & 0.252 & 1.000 & 1.000 & 1.000 \\
mistral & A & dog, guitar & aggregate & 600 & +0.0033 & 0.843 & 0.921 & 1.000 & 1.000 & 1.000 \\
mistral & A & tea, history & aggregate & 600 & -0.0183 & 0.300 & 0.341 & 1.000 & 1.000 & 1.000 \\
mistral & A & cycling, bread & near ceiling & 263 & +0.0228 & 0.273 & 0.362 & 1.000 & 1.000 & 1.000 \\
mistral & A & cycling, bread & middle & 98 & +0.0510 & 0.531 & 0.592 & 1.000 & 1.000 & 1.000 \\
mistral & A & cycling, bread & far from ceiling & 239 & +0.0084 & 0.637 & 0.819 & 1.000 & 1.000 & 1.000 \\
mistral & A & dog, guitar & near ceiling & 249 & +0.0120 & 0.491 & 0.650 & 1.000 & 1.000 & 1.000 \\
mistral & A & tea, history & near ceiling & 249 & +0.0201 & 0.225 & 0.330 & 1.000 & 1.000 & 1.000 \\
mistral & A & dog, guitar & middle & 102 & +0.0196 & 0.811 & 0.893 & 1.000 & 1.000 & 1.000 \\
mistral & A & tea, history & middle & 102 & -0.0588 & 0.502 & 0.568 & 1.000 & 1.000 & 1.000 \\
mistral & A & dog, guitar & far from ceiling & 249 & -0.0120 & 0.439 & 0.607 & 1.000 & 1.000 & 1.000 \\
mistral & A & tea, history & far from ceiling & 249 & -0.0402 & 0.012 & 0.024 & 0.149 & 0.286 & 0.298 \\
mistral & A2 & cycling, bread & aggregate & 600 & +0.1474 & 0.607 & 0.129 & 1.000 & 0.450 & 0.900 \\
mistral & A2 & dog, guitar & aggregate & 600 & +0.3698 & $<$0.001 & $<$0.001 & $<$0.001 & 0.001 & $<$0.001\,$\star$ \\
mistral & A2 & tea, history & aggregate & 600 & +0.2782 & $<$0.001 & 0.011 & 0.004 & 0.063 & 0.008\,$\star$ \\
mistral & A2 & cycling, bread & near ceiling & 52 & +0.0162 & $<$0.001 & $<$0.001 & $<$0.001 & 0.007 & 0.002\,$\star$ \\
mistral & A2 & dog, guitar & near ceiling & 52 & +0.0195 & $<$0.001 & 0.005 & $<$0.001 & 0.033 & $<$0.001\,$\star$ \\
mistral & A2 & tea, history & near ceiling & 52 & +0.1377 & $<$0.001 & $<$0.001 & $<$0.001 & 0.001 & $<$0.001\,$\star$ \\
mistral & A2 & cycling, bread & middle & 191 & +0.0944 & 0.316 & 0.098 & 1.000 & 0.450 & 0.900 \\
mistral & A2 & dog, guitar & middle & 191 & +0.1080 & 0.714 & 0.090 & 1.000 & 0.450 & 0.900 \\
mistral & A2 & tea, history & middle & 191 & +0.3261 & $<$0.001 & $<$0.001 & 0.005 & 0.002 & 0.004\,$\star$ \\
mistral & A2 & cycling, bread & far from ceiling & 357 & +0.1949 & 0.371 & 0.223 & 1.000 & 0.450 & 0.900 \\
mistral & A2 & dog, guitar & far from ceiling & 357 & +0.5608 & $<$0.001 & $<$0.001 & $<$0.001 & 0.002 & $<$0.001\,$\star$ \\
mistral & A2 & tea, history & far from ceiling & 357 & +0.2729 & 0.075 & 0.124 & 0.376 & 0.450 & 0.752 \\
mistral & B & cycling, bread & aggregate & 100 & +0.4278 & $<$0.001 & $<$0.001 & $<$0.001 & $<$0.001 & $<$0.001\,$\star$ \\
mistral & B & dog, guitar & aggregate & 100 & +0.5859 & $<$0.001 & $<$0.001 & $<$0.001 & $<$0.001 & $<$0.001\,$\star$ \\
mistral & B & tea, history & aggregate & 100 & +0.2515 & $<$0.001 & $<$0.001 & $<$0.001 & $<$0.001 & $<$0.001\,$\star$ \\
mistral & C2 & cycling, bread & aggregate, vs a & 15 & +2.6417 & $<$0.001 & $<$0.001 & $<$0.001 & $<$0.001 & $<$0.001\,$\star$ \\
mistral & C2 & dog, guitar & aggregate, vs a & 33 & +3.3366 & $<$0.001 & $<$0.001 & $<$0.001 & $<$0.001 & $<$0.001\,$\star$ \\
mistral & C2 & dog, guitar & aggregate, vs b & 12 & +1.2018 & $<$0.001 & $<$0.001 & $<$0.001 & $<$0.001 & $<$0.001\,$\star$ \\
mistral & C2 & tea, history & aggregate, vs a & 127 & +4.0434 & $<$0.001 & $<$0.001 & $<$0.001 & $<$0.001 & $<$0.001\,$\star$ \\
mistral & C2 & tea, history & aggregate, vs b & 29 & +0.7942 & $<$0.001 & $<$0.001 & $<$0.001 & $<$0.001 & $<$0.001\,$\star$ \\
qwen & A & cycling, bread & aggregate & 600 & +0.0000 & 1.000 & 1.000 & 1.000 & 1.000 & 1.000 \\
qwen & A & dog, guitar & aggregate & 600 & +0.0050 & 0.750 & 0.829 & 1.000 & 1.000 & 1.000 \\
qwen & A & tea, history & aggregate & 600 & -0.0150 & 0.361 & 0.425 & 1.000 & 1.000 & 1.000 \\
qwen & A & cycling, bread & near ceiling & 304 & -0.0164 & 0.180 & 0.183 & 1.000 & 1.000 & 1.000 \\
qwen & A & cycling, bread & middle & 86 & +0.1047 & 0.160 & 0.209 & 1.000 & 1.000 & 1.000 \\
qwen & A & cycling, bread & far from ceiling & 210 & -0.0190 & 0.317 & 0.456 & 1.000 & 1.000 & 1.000 \\
qwen & A & dog, guitar & near ceiling & 291 & -0.0241 & 0.065 & 0.063 & 0.720 & 0.693 & 1.000 \\
qwen & A & tea, history & near ceiling & 291 & +0.0137 & 0.346 & 0.469 & 1.000 & 1.000 & 1.000 \\
qwen & A & dog, guitar & middle & 89 & +0.0562 & 0.520 & 0.602 & 1.000 & 1.000 & 1.000 \\
qwen & A & tea, history & middle & 89 & -0.1798 & 0.044 & 0.054 & 0.529 & 0.648 & 1.000 \\
qwen & A & dog, guitar & far from ceiling & 220 & +0.0227 & 0.225 & 0.332 & 1.000 & 1.000 & 1.000 \\
qwen & A & tea, history & far from ceiling & 220 & +0.0136 & 0.467 & 0.634 & 1.000 & 1.000 & 1.000 \\
qwen & A2 & cycling, bread & aggregate & 600 & +0.1818 & 0.867 & 0.038 & 1.000 & 0.378 & 0.756 \\
qwen & A2 & dog, guitar & aggregate & 600 & +0.1331 & 0.483 & 0.121 & 1.000 & 0.848 & 1.000 \\
qwen & A2 & tea, history & aggregate & 600 & -0.0150 & 0.731 & 0.859 & 1.000 & 1.000 & 1.000 \\
qwen & A2 & cycling, bread & near ceiling & 264 & -0.0281 & 0.130 & 0.002 & 1.000 & 0.019 & 0.038\,$\star$ \\
qwen & A2 & dog, guitar & near ceiling & 264 & -0.0217 & 0.908 & 0.060 & 1.000 & 0.478 & 0.955 \\
qwen & A2 & tea, history & near ceiling & 264 & -0.0173 & 0.205 & 0.457 & 1.000 & 1.000 & 1.000 \\
qwen & A2 & cycling, bread & middle & 86 & -0.1237 & 0.407 & 0.424 & 1.000 & 1.000 & 1.000 \\
qwen & A2 & dog, guitar & middle & 86 & -0.1576 & 0.100 & 0.409 & 1.000 & 1.000 & 1.000 \\
qwen & A2 & tea, history & middle & 86 & -0.0600 & 0.571 & 0.772 & 1.000 & 1.000 & 1.000 \\
qwen & A2 & cycling, bread & far from ceiling & 250 & +0.5087 & 0.060 & 0.013 & 0.656 & 0.140 & 0.279 \\
qwen & A2 & dog, guitar & far from ceiling & 250 & +0.3965 & 0.052 & 0.040 & 0.618 & 0.378 & 0.756 \\
qwen & A2 & tea, history & far from ceiling & 250 & +0.0029 & 0.633 & 0.988 & 1.000 & 1.000 & 1.000 \\
qwen & B & cycling, bread & aggregate & 100 & +0.5533 & $<$0.001 & $<$0.001 & $<$0.001 & $<$0.001 & $<$0.001\,$\star$ \\
qwen & B & dog, guitar & aggregate & 100 & +0.3794 & $<$0.001 & $<$0.001 & $<$0.001 & $<$0.001 & $<$0.001\,$\star$ \\
qwen & B & tea, history & aggregate & 100 & +0.4640 & $<$0.001 & $<$0.001 & $<$0.001 & $<$0.001 & $<$0.001\,$\star$ \\
qwen & C2 & cycling, bread & aggregate, vs a & 200 & +0.6469 & $<$0.001 & $<$0.001 & $<$0.001 & $<$0.001 & $<$0.001\,$\star$ \\
qwen & C2 & cycling, bread & aggregate, vs b & 200 & +6.4400 & $<$0.001 & $<$0.001 & $<$0.001 & $<$0.001 & $<$0.001\,$\star$ \\
qwen & C2 & dog, guitar & aggregate, vs a & 200 & +2.6731 & $<$0.001 & $<$0.001 & $<$0.001 & $<$0.001 & $<$0.001\,$\star$ \\
qwen & C2 & dog, guitar & aggregate, vs b & 200 & +5.2081 & $<$0.001 & $<$0.001 & $<$0.001 & $<$0.001 & $<$0.001\,$\star$ \\
qwen & C2 & tea, history & aggregate, vs a & 200 & +4.8869 & $<$0.001 & $<$0.001 & $<$0.001 & $<$0.001 & $<$0.001\,$\star$ \\
qwen & C2 & tea, history & aggregate, vs b & 200 & +5.4903 & $<$0.001 & $<$0.001 & $<$0.001 & $<$0.001 & $<$0.001\,$\star$ \\
\end{longtable}
\end{center}
A star marks a floor that moves: decision $p$ below 0.05.



\begin{thebibliography}{99}

\bibitem{kirk2021} H. Kirk, Y. Jun, H. Iqbal, et al. \emph{Bias Out-of-the-Box:
An Empirical Analysis of Intersectional Occupational Biases in Popular
Generative Language Models.} Advances in Neural Information Processing Systems
34, 2021. arXiv:2102.04130. Code: \url{https://github.com/oxai/intersectional_gpt2}.

\bibitem{wan2025} Y. Wan, K.-W. Chang. \emph{White Men Lead, Black Women Help?
Uncovering Gender, Racial, and Intersectional Bias in Language Agency.}
Proceedings of the 63rd Annual Meeting of the Association for Computational
Linguistics, 2025, 2025.acl-long.445. Code:
\url{https://github.com/elainew728/labe-agency}.

\bibitem{wilson2024} K. Wilson, A. Caliskan. \emph{Gender, Race, and
Intersectional Bias in Resume Screening via Language Model Retrieval.}
Proceedings of the AAAI/ACM Conference on AI, Ethics, and Society, 2024.
arXiv:2407.20371. Code: \url{https://github.com/kyrawilson/Resume-Screening-Bias}.

\bibitem{beck2024} T. Beck, H. Schuff, A. Lauscher, I. Gurevych.
\emph{Sensitivity, Performance, Robustness: Deconstructing the Effect of
Sociodemographic Prompting.} Proceedings of the 18th Conference of the European
Chapter of the Association for Computational Linguistics, 2024.
arXiv:2309.07034. Code:
\url{https://github.com/UKPLab/arxiv2023-sociodemographic-prompting}.

\bibitem{ghanbari2025} S. Ghanbari Haez, M. Dragoni. \emph{Neutral Is Not
Unbiased: Evaluating Sociodemographic Shifts in Language Model Responses.}
Findings of the Association for Computational Linguistics: EMNLP 2025,
2025.findings-emnlp.814. Code:
\url{https://github.com/Saba-Gh-H/Neutral_isnot_unbiased}.

\bibitem{p1} M. Podgortsev. \emph{The Average-Fit Loss: A Unifying Metric for
the Cost of One-Size-Fits-All Systems.} 2026. Concept DOI
10.5281/zenodo.22364499.

\bibitem{p2} M. Podgortsev. \emph{Single Human-AI Fit: One Attribute, Seven
Channels, Two Directions.} 2026. Concept DOI 10.5281/zenodo.22364497.

\bibitem{kearns2017} M. Kearns, S. Neel, A. Roth, Z. S. Wu. \emph{Preventing
Fairness Gerrymandering: Auditing and Learning for Subgroup Fairness.}
arXiv:1711.05144; ICML 2018.

\bibitem{guo2020} W. Guo, A. Caliskan. \emph{Detecting Emergent Intersectional
Biases: Contextualized Word Embeddings Contain a Distribution of Human-like
Biases.} arXiv:2006.03955; AIES 2021.

\bibitem{lalor2022} J. P. Lalor, Y. Yang, K. Smith, N. Forsgren, A. Abbasi.
\emph{Benchmarking Intersectional Biases in NLP.} NAACL 2022, 3598--3609.

\bibitem{gohar2023} U. Gohar, L. Cheng. \emph{A Survey on Intersectional
Fairness in Machine Learning: Notions, Mitigation, and Challenges.}
arXiv:2305.06969; IJCAI 2023.

\bibitem{winoidentity2025} F. Arif Khan, K. Sivakumar, J. Wang, et al.
\emph{Investigating Intersectional Bias in Large Language Models using
Confidence Disparities in Coreference Resolution.} arXiv:2508.07111; COLM 2025.

\bibitem{reddy2025} S. Reddy, S. Srinivasan, U. Kanewala. \emph{Metamorphic
Testing for Fairness Evaluation in Large Language Models: Identifying
Intersectional Bias in LLaMA and GPT.} arXiv:2504.07982.

\bibitem{gupta2026} A. Gupta, P. Patil, S. Ghosh, R. Gaikwad. \emph{Compounding
Disadvantage: Auditing Intersectional Bias in LLM-Generated Explanations Across
Indian and American STEM Education.} arXiv:2601.14506.

\bibitem{ableist2025} M. Phutane, H. Jung, S. Kim, T. Mitra, A. Vashistha.
\emph{ABLEIST: Intersectional Disability Bias in LLM-Generated Hiring
Scenarios.} arXiv:2510.10998.

\bibitem{boufaied2026} A. Boufaied, D. De Souza Santos, A. Barcomb.
\emph{Intersectional Fairness in Large Language Models.} arXiv:2604.20677.

\bibitem{maltbie2026} N. Maltbie, A. Raval. \emph{Intersectional Sycophancy: How
Perceived User Demographics Shape False Validation in Large Language Models.}
arXiv:2604.11609.

\bibitem{fu2026} Y. Fu, C. Meng, X. Chen, Y. Huang, T. Wu. \emph{How Identity
and Opinion Shape Political Sycophancy in LLMs.} arXiv:2608.29198; EMNLP 2026.

\bibitem{kamruzzaman2026} M. Kamruzzaman, A. Das, G. L. Kim. \emph{Demographic
Prompting at Scale: When More Attributes Hurt LLM-Human Agreement.}
arXiv:2607.10590.

\bibitem{occhipinti2026} D. Occhipinti, M. Piergentili, M. Guerini.
\emph{LLM-as-a-Demographic: Whom Sociodemographic Prompting Helps, and Whom It
Hurts.} arXiv:2609.00222.

\bibitem{liu2026} Y. Liu, Z. Yang, P. Sheth, et al. \emph{Individual and
Combined Effects of English as a Second Language and Typos on LLM Performance.}
arXiv:2604.04723.

\bibitem{haq2026} S. Haq, B. Saldías. \emph{Dialect vs Demographics:
Quantifying LLM Bias from Implicit Linguistic Signals vs. Explicit User
Profiles.} arXiv:2604.21152; FAccT 2026.

\bibitem{tonneau2026} M. Tonneau, S. Sehgal, M. Malhotra, et al.
\emph{Different Demographic Cues Yield Inconsistent Conclusions About LLM
Personalization and Bias.} arXiv:2601.18486; EMNLP 2026.

\bibitem{guptapersona2023} S. Gupta, V. Shrivastava, A. Deshpande, et al.
\emph{Bias Runs Deep: Implicit Reasoning Biases in Persona-Assigned LLMs.}
arXiv:2311.04892; ICLR 2024.

\bibitem{poole2024} M. Poole-Dayan, D. Roy, J. Kabbara. \emph{LLM Targeted
Underperformance Disproportionately Impacts Vulnerable Users.} arXiv:2406.17737;
AAAI 2026.

\bibitem{kamruzzaman2025} M. Kamruzzaman, G. L. Kim. \emph{The Impact of
Disability Disclosure on Fairness and Bias in LLM-Driven Candidate Selection.}
arXiv:2506.00256; FLAIRS 2025.

\bibitem{panda2025} S. Panda, A. Agarwal, P. Patel. \emph{AccessEval:
Benchmarking Disability Bias in Large Language Models.} arXiv:2509.22703.

\bibitem{wang2024a} X. Wang, B. Ma, C. Hu, L. Weber-Genzel, P. Röttger, F.
Kreuter, D. Hovy, B. Plank. \emph{``My Answer is C'': First-Token Probabilities
Do Not Match Text Answers in Instruction-Tuned Language Models.}
arXiv:2402.14499; Findings of ACL 2024.

\bibitem{wang2024b} X. Wang, C. Hu, B. Ma, P. Röttger, B. Plank. \emph{Look at
the Text: Instruction-Tuned Language Models are More Robust Multiple Choice
Selectors than You Think.} arXiv:2404.08382; COLM 2024.

\bibitem{yin2025} H. Yin, M. Y. Vardi, V. Choudhary. \emph{Fragile Preferences:
A Deep Dive Into Order Effects in Large Language Models.} arXiv:2506.14092; PNAS
Nexus.

\bibitem{yildirim2026} S. Yildirim. \emph{Differential Harm Propensity in
Personalized LLM Agents: The Curious Case of Mental Health Disclosure.}
arXiv:2603.16734.

\bibitem{pashley2025} N. E. Pashley, B. Libgober, T. Dasgupta. \emph{Analysis
and sample-size determination for $2^K$ audit experiments with binary response.}
arXiv:2503.10591.

\bibitem{hoffmann2026} M. Hoffmann, C. Jouffroy, J. Jouanneau, M. Palyart, P.
Pebereau. \emph{Evaluating LLM Behavior in Hiring: Implicit Weights, Fairness
Across Groups, and Alignment with Human Preferences.} arXiv:2601.11379.

\bibitem{leyva2026} M. Leyva-Vazquez, F. Smarandache. \emph{Institutional
Prestige as Geographic Bias in Large Language Models: Evidence from Three
Factorial Experiments with Bootstrap Confidence Intervals.} arXiv:2608.18107.

\bibitem{vohra2026} S. Vohra, M. Ravikiran. \emph{The Audit Decides the Verdict:
Instrument Effects Rival Demographic Bias in LLM Decision Audits.} arXiv:2609.09048.

\bibitem{rennard2026} V. Rennard, C. Xypolopoulos. \emph{Large language models
simulate intersectional synthetic identities with a budget of one to two
dimensions.} arXiv:2608.23005.

\bibitem{magee2021} L. Magee, L. Ghahremanlou, K. Soldatic, S. Robertson.
\emph{Intersectional Bias in Causal Language Models.} arXiv:2107.07691.

\end{thebibliography}

\end{document}
